\subsection{20.89/19.83: almost Engel closure under additional hypotheses}
\label{partial:engel}

An element \(a\) is almost Engel in \(G\) if one finite set \(E(a)\)
contains \([x,{}_n a]\) for every sufficiently large \(n\), for
each \(x\in G\). The entry time may depend on \(x\).
The question asks whether these elements form a subgroup of an
arbitrary linear group. A finite sink inside
\(\langle a_1,\ldots,a_r\rangle\) does not by itself provide a
sink for a product of the \(a_i\) relative to all of \(G\).
The following deductions supply that ambient quantifier in
specified cases; the general positive-characteristic case remains
unresolved.

\paragraph{The local reduction.}
In every characteristic, a subgroup generated by finitely many
elements almost Engel in \(G\) is finite-by-nilpotent.
Here is the reduction and its imported inputs.
Restrict to the finitely generated linear group
\(X=\langle a_1,\ldots,a_r\rangle\), which is residually finite.
Choose a finite-index normal subgroup \(H\) avoiding the
nonidentity members of their finite sinks. Then each \(a_i\)
is Engel in \(H\langle a_i\rangle\): for \(h\in H\) eventual
commutators lie both in \(H\) and in its avoided sink, and
\([ha_i^k,{}_n a_i]=[h,{}_n a_i]^{a_i^k}\).
Gruenberg's theorem identifies the Engel elements of a linear
group with its Hirsch--Plotkin radical. Hence
\([H,a_i]\subseteq R=\operatorname{HP}(H)\).
In \(X/R\), all \(a_i\) centralize \(H/R\); its centralizer has
a central subgroup of finite index, namely its intersection
with \(H/R\). The locally nilpotent linear group \(R\) is soluble.
Consequently \(X\) is virtually soluble.
Shumyatsky's Theorem~3.3 now makes \(X\) finite-by-nilpotent
\cite{shumyatsky-engel}; the residual-finiteness and
linear-radical inputs are recorded also in
\cite{wehrfritz-engel}.
No product closure has been assumed in this deduction.

Let \(S\) be the set in question and \(L=\langle S\rangle\).
Every finitely generated subgroup of \(L\) is finite-by-nilpotent,
by collecting the finitely many members of \(S\) used in its words.
Call this property \(\mathcal P\).
In any group with \(\mathcal P\), the torsion elements form a
locally finite characteristic subgroup: finitely many torsion
elements generate a finite-by-nilpotent group whose nilpotent
quotient is generated by finitely many torsion elements, hence
finite. This proves closure as well as local finiteness.

The unique minimal finite sink consists of the periodic points
of \(x\mapsto[x,a]\). Each such point belongs to
\(\langle a^G\rangle\). Express its finitely many points using
finitely many conjugates of \(a\), including \(a\) itself.
Their subgroup is finite-by-nilpotent by the local reduction.
In its nilpotent quotient a periodic repeated commutator must
be trivial, so the entire minimal sink lies in the finite kernel.
Thus its generated subgroup is finite. It need not be normal
in the ambient group. In particular, for a torsion-free linear
group every almost Engel element is Engel, proving closure there
in every characteristic.

\paragraph{Two lifting facts.}
If \(X/D\) is locally nilpotent for some finite normal \(D\),
one can choose \(D\) characteristic: take one of minimum order.
For an automorphism \(\sigma\), the quotient by
\(D\cap\sigma(D)\) embeds in the product of the two locally
nilpotent quotients. That product is locally nilpotent, as are
its subgroups. Minimality forces \(D=\sigma(D)\).

Suppose \(X\) has \(\mathcal P\), \(N\lhd X\) is torsion-free,
and \(X/N\) is finite-by-(locally nilpotent).
Choose a finite characteristic kernel in \(X/N\) and let \(R\)
be its preimage. The torsion subgroup \(K=T(R)\) is locally
finite and injects into the finite group \(R/N\), so is finite.
For a finitely generated \(Y\le X\), a sufficiently late lower
central term is finite by \(\mathcal P\), and a sufficiently
late term lies in \(R\) since \(X/R\) is locally nilpotent.
A common later term therefore lies in \(K\).
Thus \(X/K\) is locally nilpotent. The construction preserves
any group of automorphisms preserving \(N\).

\paragraph{Characteristic zero.}
The group \(L\) contains no nonabelian free subgroup, since such
a group cannot have \(\mathcal P\).
The characteristic-zero form of the Tits alternative applies
without finite generation and makes \(L\) virtually soluble
\cite[Theorem~1]{tits}.
Over an algebraic closure of the field let \(A\) be its Zariski
closure and \(U\) the unipotent radical of \(A^\circ\).
The connected component is soluble, and \(A^\circ/U\) is a torus;
the quotient \(A/U\) is linear algebraic
\cite{milne-algebraic}.
Since \(L\lhd G\), the subgroup \(N=L\cap U\) is normal in \(G\).
In characteristic zero it is torsion-free: the finite matrix
logarithm gives \(\log(u^m)=m\log(u)\).
The group \(L/N\) is therefore linear and virtually abelian.

Shumyatsky's Lemma~3.2 gives the following prior consequence:
in a virtually abelian group, almost Engel elements form a
subgroup. Indeed it makes the normal closure of each such element
finite-by-(locally nilpotent); choose its finite kernel characteristic
as above and use the Hirsch--Plotkin radical after factoring out
the finitely many kernels needed for a product.
The images of \(S\) generate \(L/N\), so every element of \(L/N\)
is almost Engel there. Only now apply Shumyatsky's Theorem~4.2:
an almost Engel linear group is finite-by-hypercentral,
hence finite-by-(locally nilpotent).
The torsion-free lifting fact gives a finite \(G\)-normal
subgroup \(K\le L\) with \(L/K\) locally nilpotent.
For \(l\in L,x\in G\), normality gives \([x,l]\in L\).
The two elements \([x,l]K,lK\) generate a nilpotent subgroup,
so \([x,{}_n l]\in K\) eventually. Thus \(K\) is one ambient
sink for every \(l\), proving \(S=L\).
The minimum-kernel argument makes \(K\) characteristic in \(L\),
and hence in \(G\).

This proves closure for every characteristic-zero linear group,
without generation or solubility assumptions on \(G\).
It does not make the full torsion subgroup of \(L\) finite;
an infinite central torsion subgroup is possible.
Both the unrestricted Tits step and the torsion-free unipotent
kernel require characteristic zero in this argument.

\paragraph{Two further characteristic-free cases.}
If all finite subgroups of \(G\) have order at most \(b\),
the locally finite torsion subgroup \(K=T(L)\) has order at
most \(b\): \(b+1\) of its elements would generate a larger
finite subgroup. Property \(\mathcal P\) then gives \(L/K\)
locally nilpotent, and the same commutator argument proves
\(S=L\), with \(K\) a characteristic common sink.
If \(G\) has a torsion-free subgroup of index \(d\), every
finite subgroup acts freely on its \(d\) cosets, so its order
divides \(d\); in particular \(|K|\mid d\).
This is not automatic in positive characteristic: the group
generated by \(\operatorname{diag}(t,1)\) and
\(\left(\begin{smallmatrix}1&1\\0&1\end{smallmatrix}\right)\)
over \(\F_p(t)\) contains elementary abelian subgroups of
orders \(p^r\) for all \(r\).

For an upper triangular group, order matrix entries by decreasing
superdiagonal distance and refine its unipotent radical into
one-entry additive factors \(B_s\). Conjugation acts on the
factor at \((i,j)\) by the character
\(\chi_s(g)=g_{jj}/g_{ii}\). Define
\[
 P=\bigcap_{B_s\text{ infinite}}\ker\chi_s.
\]
If multiplication by \(\chi-1\ne0\) on an additive subgroup of
a field has one finite eventual sink, a nonzero orbit first
shows that \(\chi-1\) has finite multiplicative order.
All points are then periodic and must lie in the sink, so the
additive subgroup is finite. This proves necessity of membership
in \(P\).
Conversely \(P\) has a finite normal series whose factors are
finite or \(P\)-central, including its abelian diagonal quotient.
The Hall--Baer finite-central-series lemma makes it
finite-by-nilpotent \cite{wehrfritz-engel}.
A finite lower-central term \(K=\gamma_{c+1}(P)\) is normal
in \(G\); since \([G,P]\subseteq P\),
\([x,{}_{c+1}g]\in K\) for \(x\in G,g\in P\).
This proves that \(P\) is exactly the almost Engel subgroup,
with a uniform commutator-length bound in this triangular case.

\paragraph{Normal subgroups of algebra unit groups.}
Let \(A\) be a finite-dimensional associative unital algebra over
an infinite field \(F\), and \(G\lhd A^\times\).
Every almost Engel element of \(G\) is Engel.
For \(a\in G\), all but finitely many \(c_t=1+ta\) are units:
\(\det(I+tL_a)\), in the faithful left regular representation,
is a nonzero polynomial with constant term one.
Each such \(c_t\) commutes with \(a\) and normalizes \(G\), hence
permutes its minimal sink \(E(a)\).
If \(c_s^{-1}ec_s=c_t^{-1}ec_t=f\) for \(s\ne t\), then
\[
 e+s ea=f+s af,\qquad e+t ea=f+t af.
\]
Subtracting gives \(ea=af\), and substitution gives \(e=f\);
thus \(e\) commutes with \(a\).
Finiteness of \(E(a)\) and infinitely many parameters give such
a repetition for every \(e\). The commutator map consequently
sends the whole sink to \(1\); since it permutes the minimal
sink, that sink is \(\{1\}\). Gruenberg's theorem now proves
the subgroup assertion. This includes normal subgroups of
\(\GL_n(F)\) and \(\SL_n(F)\).
Over \(\F_q\), the same argument gives
\(|E(a)|\ge q-\dim_F A+1\) for a nontrivial minimal sink;
in a matrix subalgebra it sharpens to \(q-s(a)+1\), where
\(s(a)\) counts the distinct eigenvalues in \(\F_q\).

Normality supplies the parameter conjugations. An arbitrary
finite \(S_3\le\GL_2(\Q)\) has non-Engel almost Engel elements,
so this hypothesis cannot simply be dropped.
Proofs, exact finite controls, and the precise classical imports
are indexed by \artifact{research/20.89-char0-proof.md},
\artifact{research/20.89-residual-reduction.md},
\artifact{research/20.89-triangular-proof.md},
\artifact{research/20.89-bounded-torsion-proof.md}, and
\artifact{research/20.89-normal-algebra-proof.md}.
